\section{Illustrative Fits under Alternative Mating Models}
\label{app:assortment-comparison}

Changing the mating mechanism changes both the predicted correlations and
what the data can identify. Using Clark's notation, let $m$ be the spouse
genetic correlation, $r$ the spouse correlation in the true phenotype,
$h^2$ its heritability, and $\theta$ the common attenuation factor. Under
phenotypic assortment, $m=h^2r$. An alternative in which spouses assort
only on additive genetic value, with environmental components having
conditional mean zero given own genetic value, instead gives $r=h^2m$.
Retaining additive transmission, fresh offspring environments, and the
corresponding conditionally independent outside-mate construction, both
models predict sibling correlation $\theta h^2(1+m)/2$. But phenotypic
assortment predicts parent--child correlation $\theta(h^2+m)/2$, whereas
genetic assortment makes it equal to the sibling correlation. Likewise,
genetic assortment equates grandparent and avuncular correlations.
The common shape of the collateral-relative formulas therefore does not
justify fitting the original table with three unrestricted structural
parameters and interpreting it as either model.

To illustrate, we fitted both models by unweighted nonlinear least squares
to the eleven published higher-education correlations for births
1860--1919 in \citet[Table~2]{Clark2023}. These are the rounded PNAS
correlations, not the different correlations in the book's Table~4.5.
The fits impose the appropriate parameter bounds and assortment
restriction; no observed spouse correlation is included.

\begin{table}[htbp]
\centering
\caption{Illustrative constrained fits to PNAS education correlations}
\label{tab:assortment-comparison}
\begin{tabular}{lrrrrr}
\hline
Model & $\theta$ & $h^2$ & $m$ & Implied $r$ & SSE\\
\hline
Genetic, no error & 1 & 0.4278 & 0.5638 & 0.2412 & 0.01029\\
Phenotypic, no error & 1 & 0.4393 & 0.4393 & 1.0000 & 0.02701\\
Phenotypic, error allowed & 0.5176 & 0.7765 & 0.5810 & 0.7482 & 0.00928\\
\hline
\end{tabular}
\end{table}

Without measurement error, the phenotypic model reaches the boundary of
perfect spouse phenotype correlation. Allowing error yields an admissible
fit and a smaller residual sum of squares, with an additional identified
parameter. Under genetic assortment, by contrast, every moment depends on
$\theta$ and $h^2$ only through their product: freeing $\theta$ leaves the
fit unchanged, with $\theta h^2=0.4278$ and $m=0.5638$. Even an observed
spouse correlation would identify only $\theta h^2m$. Under phenotypic
assortment, the lineal--collateral contrast can instead supply the
information needed to separate heritability and attenuation, as derived
in Section~\ref{sec:me-identification}.

These calculations sharpen rather than settle the criticism. A high
genetic spouse correlation relative to the observed genetic variance share
can be consistent with a properly specified latent-phenotype model.
The requirement is to derive, impose, and check that model's restrictions
throughout, not to invoke a different matching mechanism while retaining
the original parameter interpretations---the issue highlighted by
\citet[p.~19]{Goldberger1978}. The fits here are descriptive, not formal
specification tests: correlations have dependent sampling errors, and our
equal-sex measurement and heritability assumptions are particularly strong
for historical education, which Clark records only for men. The fitting
code and numerical output are supplied in the repository's
\texttt{verification/} directory.

The appendix's stated estimation procedure warrants a precise distinction.
Equation A3 specifies three freely fitted regression coefficients; it
does not state a nonlinear constraint or a constrained recovery of the
structural parameters. This has one coefficient too many for phenotypic
assortment with known attenuation, or for genetic assortment, whose lineal
adjustment must vanish. With unknown attenuation under phenotypic
assortment, however, three coefficients can be legitimate: recover $m$
and $r$, impose $h^2=m/r$, and recover $\theta$ from the intercept,
checking admissibility. The appendix does not explain that recovery;
its intercept identifies $\theta h^2$, not latent heritability alone.
The number of coefficients therefore establishes the excess freedom in
the first two interpretations, but does not by itself establish it in
the third.
